\documentclass[11pt,a4paper]{article}

% ------------------------------------------------------------ fonts / layout
\usepackage[T1]{fontenc}
\usepackage[utf8]{inputenc}
\usepackage{lmodern}
\usepackage{microtype}
\usepackage[margin=1in]{geometry}
\usepackage{xcolor}
\usepackage{graphicx}
\usepackage{booktabs}
\usepackage{amsmath}
\usepackage{float}
\usepackage{enumitem}
\usepackage{fancyhdr}
\usepackage{titlesec}
\usepackage{caption}
\usepackage{needspace}
\usepackage{etoolbox}
\usepackage{listings}
\usepackage[hidelinks]{hyperref}

\graphicspath{{../results/}}

% ------------------------------------------------------------------- colours
\definecolor{accent}{HTML}{1F4E79}      % headings, rules
\definecolor{codebg}{HTML}{F6F8FA}      % listing background
\definecolor{codeframe}{HTML}{D0D7DE}
\definecolor{codekw}{HTML}{0550AE}      % keywords
\definecolor{codecom}{HTML}{6E7781}     % comments
\definecolor{codenum}{HTML}{8C959F}     % line numbers
\definecolor{codestr}{HTML}{0A3069}

% ------------------------------------------------------- paragraphs / lists
\linespread{1.08}
\setlength{\parindent}{0pt}
\setlength{\parskip}{0.55em}
\setlist[itemize]{leftmargin=1.5em, itemsep=0.25em, topsep=0.3em, parsep=0pt}
\setlist[enumerate]{leftmargin=1.7em, itemsep=0.25em, topsep=0.3em, parsep=0pt}

% ----------------------------------------------------------------- headings
\titleformat{\section}
    {\Large\bfseries\color{accent}}{\thesection}{0.7em}{}[\vspace{-0.3em}\color{accent}\titlerule]
\titleformat{\subsection}
    {\large\bfseries\color{accent}}{\thesubsection}{0.7em}{}
\titleformat{\subsubsection}
    {\normalsize\bfseries}{\thesubsubsection}{0.7em}{}
\titleformat{\paragraph}[hang]
    {\normalfont\bfseries\color{accent}}{}{0em}{}
\titlespacing*{\section}{0pt}{1.6em}{0.9em}
\titlespacing*{\subsection}{0pt}{1.3em}{0.5em}
\titlespacing*{\subsubsection}{0pt}{1.0em}{0.4em}
\titlespacing*{\paragraph}{0pt}{1.1em}{0.4em}

% keep a heading together with at least a few lines of the text that follows it
\pretocmd{\section}{\needspace{8\baselineskip}}{}{}
\pretocmd{\subsection}{\needspace{6\baselineskip}}{}{}
\pretocmd{\subsubsection}{\needspace{5\baselineskip}}{}{}
\pretocmd{\paragraph}{\needspace{5\baselineskip}}{}{}

% --------------------------------------------------------- header / footer
\pagestyle{fancy}
\fancyhf{}
\fancyhead[L]{\small\color{codecom} CS683 Programming Assignment -2 report}
\fancyhead[R]{\small\color{codecom} \leftmark}
\fancyfoot[C]{\small\thepage}
\renewcommand{\headrulewidth}{0.4pt}
\renewcommand{\headrule}{\hbox to\headwidth{\color{codeframe}\leaders\hrule height \headrulewidth\hfill}}
\renewcommand{\sectionmark}[1]{\markboth{#1}{#1}}
\setlength{\headheight}{14pt}

% ------------------------------------------------------- tables / captions
\renewcommand{\arraystretch}{1.25}
\captionsetup{font=small, labelfont={bf,color=accent}, skip=6pt}

% -------------------------------------------------------------- code listings
\lstdefinestyle{champsim}{
    language=C++,
    basicstyle=\ttfamily\footnotesize,
    keywordstyle=\color{codekw}\bfseries,
    commentstyle=\color{codecom}\itshape,
    stringstyle=\color{codestr},
    numberstyle=\tiny\color{codenum},
    numbers=left, numbersep=16pt,
    backgroundcolor=\color{codebg},
    frame=l, rulecolor=\color{accent}, framerule=2pt,
    framexleftmargin=6pt, xleftmargin=3.4em,
    aboveskip=0.1em, belowskip=0pt,
    breaklines=true, breakatwhitespace=false,
    columns=fullflexible, keepspaces=true, tabsize=4,
    showstringspaces=false,
    morekeywords={uint8_t, uint16_t, uint32_t, uint64_t, int64_t, CACHE, assert},
}
\lstset{style=champsim}

% Marks a spot that still has to be filled in (red box); easy to grep for.
\newcommand{\todo}[1]{\colorbox{red!12}{\textcolor{red!70!black}{\small [TODO: #1]}}}

% ---------------------------------------------------------------- title page
\begin{document}

\begin{titlepage}
    \thispagestyle{empty}
    \centering
    \vspace*{0.28\textheight}
    {\color{accent}\rule{0.8\textwidth}{1.2pt}\par}
    \vspace{1.0cm}
    {\fontsize{21}{26}\selectfont\bfseries\color{accent} CS683 Programming Assignment -2 report\par}
    \vspace{1.0cm}
    {\color{accent}\rule{0.8\textwidth}{1.2pt}\par}
    \vspace{1.4cm}
    {\large Submitted by: 24B1015, 24B0995, 24B1021, 24B1037\par}
    \vfill
\end{titlepage}

% ================================================================= TASK 1
\section{Task 1: Stride Pattern Prefetching}

% ---------------------------------------------------------------- Task 1A
\subsection{Task 1A: Understanding the IP-stride prefetcher}
\label{sec:task1a}

\subsubsection{How the prefetcher works}

The IP-stride prefetcher (\texttt{prefetcher/ip\_stride.l1d\_pref}) sits at the L1 data
cache. Its idea is that a single load or store instruction, identified by its instruction
pointer (IP), tends to walk through memory with a constant stride (for example, a loop
over an array). It remembers the last address each IP touched, computes the stride between
consecutive accesses of that IP, and once the same stride is seen twice in a row it
prefetches the next few addresses along that stride. The code is walked through below,
section by section; the line numbers in the listings are the line numbers in the source file.

\paragraph{1.\ Parameters and the tracker entry (lines 4--30).}
\lstinputlisting[language=C++, firstline=4, lastline=30, firstnumber=4]{../prefetcher/ip_stride.l1d_pref}
\begin{itemize}
    \item \texttt{IP\_TRACKER\_COUNT = 24}: the prefetching table has only 24 entries.
    \item \texttt{PREFETCH\_DEGREE = 3}: when the prefetcher fires it issues 3 prefetches. The
          degree is fixed and never changes.
    \item \texttt{MASK\_IP = 0xFF}: only the low 8 bits of the IP are used to identify an
          instruction. Different instructions whose IPs share the low 8 bits therefore
          share one tracker (aliasing).
    \item Each \texttt{IP\_TRACKER} entry stores: the (masked) IP it belongs to; the cache-line
          address of the last access by that IP (\texttt{last\_cl\_addr}); the stride observed
          between its last two accesses (\texttt{last\_stride}, signed, in cache lines); and an
          LRU rank used to choose which entry to replace when the table is full.
    \item The table \texttt{trackers[]} is a single global array, so it is shared by everything that
          uses this prefetcher (fine for the single-core configuration used here).
\end{itemize}

\paragraph{2.\ Initialisation (lines 33--39).}
\lstinputlisting[language=C++, firstline=33, lastline=39, firstnumber=33]{../prefetcher/ip_stride.l1d_pref}
This function is called once when the simulation starts. It prints a banner and gives the 24
entries the distinct LRU ranks $0,1,\dots,23$. Rank 0 means most recently used and rank 23
(\texttt{IP\_TRACKER\_COUNT-1}) means least recently used, so at any time the ranks form a
permutation of $0\ldots23$ and exactly one entry is the replacement victim. All entries start
with \texttt{ip = 0}, \texttt{last\_cl\_addr = 0} and \texttt{last\_stride = 0}.

\paragraph{3.\ Entry point: cache-line address and IP (lines 41--46).}
\lstinputlisting[language=C++, firstline=41, lastline=46, firstnumber=41]{../prefetcher/ip_stride.l1d_pref}
\texttt{l1d\_prefetcher\_operate} is called by the L1D for each access it handles. The
prefetcher works on cache lines, not bytes: shifting the byte address right by
\texttt{LOG2\_BLOCK\_SIZE} drops the offset inside the 64-byte line, giving \texttt{cl\_addr}.
The IP is reduced to its low 8 bits. The arguments \texttt{cache\_hit} and \texttt{type} are
not used in this file, so the prefetcher trains on every access it is called for, whether it
hit or missed in the L1D.

\paragraph{4.\ Looking up the tracker (lines 48--52).}
\lstinputlisting[language=C++, firstline=48, lastline=52, firstnumber=48]{../prefetcher/ip_stride.l1d_pref}
The table is fully associative: a linear search compares the masked IP against every entry.
If an entry matches, the loop \texttt{break}s with \texttt{index} pointing at it. If none
matches, the loop runs to the end and \texttt{index} equals \texttt{IP\_TRACKER\_COUNT}.

\paragraph{5.\ A new IP: allocating a tracker (lines 54--74).}
\lstinputlisting[language=C++, firstline=54, lastline=74, firstnumber=54]{../prefetcher/ip_stride.l1d_pref}
\begin{itemize}
    \item \textbf{Choosing the victim (57--60).} The entry whose LRU rank is 23 (the least
          recently used) is selected for replacement.
    \item \textbf{Initialising it (62--64).} The entry is overwritten with the new IP and the
          current line address. \texttt{last\_stride} is set to 0 because no stride is known yet.
    \item \textbf{LRU update (67--71).} Every entry that was more recently used than the victim
          (rank smaller than the victim's) moves one step towards ``old'' (rank $+1$), and the new
          entry becomes rank 0. Since the victim had rank 23, all other entries are incremented
          and the ranks remain a permutation.
    \item \textbf{No prefetch (73).} With only one address seen, there is no stride to
          compute, so the function returns without prefetching.
\end{itemize}

\paragraph{6.\ Sanity check and stride calculation (lines 76--90).}
\lstinputlisting[language=C++, firstline=76, lastline=90, firstnumber=76]{../prefetcher/ip_stride.l1d_pref}
When execution reaches here, a tracker for this IP exists. The \texttt{assert} guards against
an impossible state. The stride is the difference between the current and the previous cache-line
address of this IP, in units of cache lines. Because the addresses are unsigned, the code
subtracts the smaller from the larger and negates the result if the access moved backwards, so a
descending access pattern gives a negative stride.

\paragraph{7.\ Ignoring repeated lines (lines 94--96).}
\lstinputlisting[language=C++, firstline=94, lastline=96, firstnumber=94]{../prefetcher/ip_stride.l1d_pref}
A stride of 0 means the same cache line was accessed again (for example several accesses to the
same line). There is nothing to learn or prefetch from that, so the function returns early. Note
that this early return does \emph{not} update \texttt{last\_cl\_addr}, \texttt{last\_stride} or the
LRU ranks.

\paragraph{8.\ The trigger and the prefetch loop (lines 98--117).}
\lstinputlisting[language=C++, firstline=98, lastline=117, firstnumber=98]{../prefetcher/ip_stride.l1d_pref}
\begin{itemize}
    \item \textbf{Trigger (100).} Prefetching only happens if the new stride equals the stored
          \texttt{last\_stride}, i.e.\ the same stride has been seen twice in a row. This is the
          prefetcher's only form of confidence: one repetition. Since a new entry starts with
          \texttt{last\_stride = 0}, an IP needs three accesses before its first prefetch (the first
          allocates the entry, the second records the stride, the third confirms it).
    \item \textbf{Addresses (103--104).} For $i = 0,1,2$ it prefetches the line at
          \texttt{cl\_addr + stride $\cdot (i+1)$}, converted back to a byte address, so with stride $s$
          it prefetches the next three lines along the pattern. 
    \item \textbf{Page boundary (108--109).} The prefetch address must be in the same 4~KB page as
          the demand access (the usual reason is that consecutive
          pages are not necessarily contiguous or even mapped, so a stride run cannot safely be
          followed across a page. Also the prefetcher does not know whether the next page is actually valid/mapped.). On the first violation the loop
          stops, since the later prefetches in the same direction would also be outside the page.
    \item \textbf{Fill level (112--115).} \texttt{prefetch\_line} hands the request to the cache. If
          fewer than half of the L1D MSHR entries are in use the line is brought into the L1D
          (\texttt{FILL\_L1}); otherwise it is only brought into the L2 (\texttt{FILL\_L2}), so that
          prefetches do not use up the miss-handling resources that demand misses need.
\end{itemize}

\paragraph{9.\ Updating the entry (lines 119--128).}
\lstinputlisting[language=C++, firstline=119, lastline=128, firstnumber=119]{../prefetcher/ip_stride.l1d_pref}
Whether or not a prefetch was issued, the entry now remembers the current line address and the
current stride, ready for the next access of this IP. The entry then becomes the most recently
used: every entry that was more recent than it is aged by one and it gets rank 0. Note that only
the \emph{last} stride is kept as history; one irregular access overwrites it.

\paragraph{10.\ Fill hook and final statistics (lines 130--139).}
\lstinputlisting[language=C++, firstline=130, lastline=139, firstnumber=130]{../prefetcher/ip_stride.l1d_pref}
\texttt{l1d\_prefetcher\_cache\_fill} is called whenever a line is filled into the L1D. This
prefetcher does not use it (the body just returns). \texttt{l1d\_prefetcher\_final\_stats} prints the
degree at the end of the simulation. There is no accuracy tracking in this prefetcher; the
accuracy reported in the results comes from ChampSim's own prefetch counters.

\subsubsection{Limitations}
The IP-stride prefetcher only works for accesses with one constant stride per IP. Its main
limitations, which motivate SPP, are:
\begin{itemize}
    \item \textbf{No tolerance to irregularity.} Each entry stores a single stride and a single
          previous address, and a prefetch is issued only when $s_n = s_{n-1}$, where
          $s_n = A_n - A_{n-1}$. One irregular access overwrites \texttt{last\_stride}, so the
          next access cannot match it and the stride has to be re-learned and re-confirmed. A single
          irregular access therefore suppresses prefetching for about two subsequent accesses.
          There is no confidence counter that could absorb such noise.
    \item \textbf{Cannot learn alternating or periodic strides.} Patterns whose stride sequence is
          periodic but not constant, such as $\{+1,+3,+1,+3,\dots\}$, never produce two equal
          consecutive strides, so no prefetch is ever issued even though the pattern is fully
          predictable. Typical sources are structure-of-fields traversals, padded 2D arrays and
          nested loops. Only one delta of history is kept, so such patterns cannot be captured.
    \item \textbf{No support for irregular or indirect accesses.} Pointer chasing and indirect
          accesses such as \texttt{A[B[i]]} have no constant stride, so the prefetcher gives no
          benefit and only trains on noise.
    \item \textbf{Training delay.} An IP needs three accesses before its first prefetch, which
          hurts short loops and short streams.
    \item \textbf{Page boundary.} Prefetching stops at the 4\,KB page boundary and each new page
          has to be trained again, so the lookahead is truncated near the end of every page.
    \item \textbf{Fixed degree and no feedback.} The degree is fixed and is not adapted to
          accuracy or timeliness, so the prefetcher may be too aggressive on irregular code
          (cache pollution, wasted bandwidth) or too shallow on regular code.
    \item \textbf{Per-IP table pressure.} Training state is kept per IP in a small LRU table.
          Programs with many load IPs thrash it, and patterns learned by one IP cannot be shared
          with other IPs that behave identically.
\end{itemize}
SPP addresses these by keeping a history of several recent deltas (a signature) and by predicting
the next delta from that history, so alternating patterns can be captured and a single irregular
access does not destroy the learned state.

\subsubsection{Results}

All runs use 25M warmup and 25M simulated instructions, with LRU at the L2C. The baseline is the
same configuration with no L1D prefetcher. Speedup is $\mathrm{IPC}_{\text{IP-stride}} /
\mathrm{IPC}_{\text{baseline}}$, and coverage is $(\text{L1D load misses}_{\text{base}} -
\text{L1D load misses}_{\text{IP-stride}}) / \text{L1D load misses}_{\text{base}}$. Accuracy is the
\texttt{ACCURACY} value printed by ChampSim for the L1D, which is the number of useful prefetches
divided by the number of prefetches \emph{issued to the lower level} (the prefetches that missed
in the L1D and went to the L2), not by all prefetches requested by the prefetcher.

\begin{table}[H]
    \centering
    \begin{tabular}{llrrrrr}
        \toprule
        Trace & Prefetcher & IPC & Speedup & L1D load MPKI & Accuracy (\%) & Coverage (\%) \\
        \midrule
        Trace1-001 & Baseline  & 0.2323 & 1.000 & 10.68 & --   & --  \\
                   & IP-stride & 0.2333 & 1.005 & 10.12 & 53.0 & 5.3 \\
        \midrule
        Trace2-002 & Baseline  & 0.2506 & 1.000 & 13.54 & --   & --  \\
                   & IP-stride & 0.2505 & 1.000 & 13.44 & 38.2 & 0.7 \\
        \midrule
        Trace3-001 & Baseline  & 0.3820 & 1.000 & 3.40  & --   & --  \\
                   & IP-stride & 0.3829 & 1.002 & 3.32  & 16.0 & 2.5 \\
        \bottomrule
    \end{tabular}
    \caption{Baseline versus IP-stride (degree 3) on the three traces.}
    \label{tab:task1a-main}
\end{table}

Table~\ref{tab:task1a-main} shows that IP-stride is almost neutral on these traces. Table~\ref{tab:task1a-pf}
gives the prefetch counters behind the accuracy and coverage numbers.

\begin{table}[H]
    \centering
    \begin{tabular}{lrrrrrr}
        \toprule
        Trace & Queued & Hit in L1D & To lower level & Useful & Useless & Late \\
        \midrule
        Trace1-001 & 155\,922 & 102\,542 & 30\,956 & 16\,408 & 8\,834  & 2\,807 \\
        Trace2-002 & 78\,701  & 57\,796  & 11\,426 & 4\,361  & 6\,300  & 272    \\
        Trace3-001 & 103\,500 & 70\,396  & 17\,007 & 2\,713  & 10\,595 & 3\,067 \\
        \bottomrule
    \end{tabular}
    \caption{L1D prefetch counters of IP-stride. ``Queued'' is the number of prefetch requests put
    into the L1D prefetch queue (the \texttt{L1D PQ ACCESS} line). ``Hit in L1D'' is the number of
    queued prefetches whose line was already in the L1D. ``To lower level'' is the denominator
    ChampSim uses for accuracy. No prefetch crosses a page boundary.}
    \label{tab:task1a-pf}
\end{table}

\paragraph{Observations.}
\begin{itemize}
    \item \textbf{Performance gain is negligible.} The speedups are $1.005$, $1.000$ and $1.002$,
          so IP-stride changes IPC by at most 0.5\%.
    \item \textbf{Coverage is low}, only $5.3\%$, $0.7\%$ and $2.5\%$ of the baseline L1D load misses
          are removed. Trace2 gains almost nothing, and it is also the trace with the highest baseline
          L1D load MPKI ($13.5$), so there is plenty of room that the prefetcher does not use.
    \item \textbf{Accuracy varies a lot between traces}: $53.0\%$, $38.2\%$ and $16.0\%$. In Trace3
          the useless prefetches ($10\,595$) are almost four times the useful ones ($2\,713$), and the
          late prefetches ($3\,067$) outnumber the useful ones, so many prefetches also arrive after
          the demand access.
    \item \textbf{Most prefetches never leave the L1D.} Of the prefetches that look up the L1D, $80.2\%$,
          $84.4\%$ and $84.1\%$ (Trace1, 2, 3) find the line already present (the \texttt{L1D PREFETCH}
          hit rate), so they fetch nothing. Only $19.9\%$, $14.5\%$ and $16.4\%$ of the queued
          prefetches go to the L2 (``To lower level'' over ``Queued''). The rest hit in the L1D or
          merge with a request already in the queue (about 8\% of the queued prefetches). The
          prefetcher therefore re-requests the same lines many times: lookups with no benefit.
    \item \textbf{ChampSim's ``ISSUED'' counts every prefetch twice.} In this build, \texttt{pf\_issued}
          is incremented once when the prefetch enters the queue and again when its translation
          request is sent to the STLB, so $\texttt{ISSUED} = \texttt{PQ ACCESS} + \texttt{PQ TO\_CACHE}$
          exactly in all three traces (for Trace1, $155\,922 + 142\,926 = 298\,848$). We therefore
          report the number of queued prefetches instead of ``ISSUED''.
    \item \textbf{Page-boundary rule works.} There are zero cross-page prefetches in all traces,
          as the code intends.
    \item \textbf{Trace3 has the least to gain}: its baseline L1D load MPKI is already low ($3.40$),
          so even a perfect prefetcher could remove only a small number of misses.
\end{itemize}
These results are the reference point for the SPP evaluation. 
% ---------------------------------------------------------------- Task 1B
\subsection{Task 1B: Stride Pattern Prefetcher (SPP)}
\label{sec:task1b}

\subsubsection{Design}
\todo{prefetching table, uniform/non-uniform classification, signature, fallback, degree, page limit}

\subsubsection{Storage overhead}
\todo{exact bit-level breakdown of every table entry and the total}

\begin{table}[H]
    \centering
    \begin{tabular}{llrrr}
        \toprule
        Table & Field & Bits per entry & Entries & Total bits \\
        \midrule
        \todo{} & & & & \\
        \bottomrule
    \end{tabular}
    \caption{Storage overhead of SPP.}
\end{table}

\subsubsection{Accuracy-based throttling}
\paragraph{Threshold selection.}
\todo{what the accuracy thresholds are and how they were decided}

\paragraph{Choice of prefetch degrees.}
\todo{why these degrees for the different accuracy ranges}

\subsubsection{Results}
\todo{IPC, speedup, L1D load MPKI, accuracy, coverage: baseline vs IP-stride vs SPP}

\begin{table}[H]
    \centering
    \begin{tabular}{llrrrrr}
        \toprule
        Trace & Prefetcher & IPC & Speedup & L1D load MPKI & Accuracy (\%) & Coverage (\%) \\
        \midrule
        \todo{} & & & & & & \\
        \bottomrule
    \end{tabular}
    \caption{Task 1 results.}
\end{table}

\begin{figure}[H]
    \centering
    \todo{speedup plot}
    \caption{Speedup over the no-prefetching baseline.}
\end{figure}

\begin{figure}[H]
    \centering
    \todo{L1D load MPKI plot}
    \caption{L1D load MPKI.}
\end{figure}

\begin{figure}[H]
    \centering
    \todo{accuracy plot}
    \caption{L1D prefetch accuracy.}
\end{figure}

\begin{figure}[H]
    \centering
    \todo{coverage plot}
    \caption{L1D prefetch coverage.}
\end{figure}

\subsubsection{Discussion}
\todo{where SPP wins or loses against IP-stride, and why}

% ================================================================= TASK 2
\section{Task 2: Cache Replacement Policy}

% ---------------------------------------------------------------- Task 2A
\subsection{Task 2A: Instruction-conscious L2 replacement (I-RRIP)}
\label{sec:task2a}

\subsubsection{Design}
\todo{changes to \texttt{update\_replacement\_state}, insertion, promotion, eviction}

\subsubsection{Rationale}
\todo{why insertion, promotion and eviction were modified or not modified, and why this choice is best}

\subsubsection{Results}
\todo{IPC, speedup, Load MPKI of L2C, L2 Instruction and L2 Data}

\begin{table}[H]
    \centering
    \begin{tabular}{llrrrrr}
        \toprule
        Trace & Policy & IPC & Speedup & L2C load MPKI & L2 instr.\ MPKI & L2 data MPKI \\
        \midrule
        \todo{} & & & & & & \\
        \bottomrule
    \end{tabular}
    \caption{Task 2A results.}
\end{table}

\todo{plots}

% ---------------------------------------------------------------- Task 2B
\subsection{Task 2B: Dynamic I-RRIP}
\label{sec:task2b}

\subsubsection{Design}
\todo{set dueling, leader sets, decision formula}

\subsubsection{Rationale for the leader sets and decision logic}
\todo{number of leader sets and why; decision logic and why}

\subsubsection{Sensitivity study}
\todo{5 values of the decision parameter (e.g.\ threshold) and their effect}

\begin{table}[H]
    \centering
    \begin{tabular}{lrrrrr}
        \toprule
        Parameter value & IPC & Speedup & L2C load MPKI & L2 instr.\ MPKI & L2 data MPKI \\
        \midrule
        \todo{} & & & & & \\
        \bottomrule
    \end{tabular}
    \caption{Sensitivity of Dynamic I-RRIP to its decision parameter.}
\end{table}

\subsubsection{Results}
\todo{IPC, speedup, Load MPKI of L2C, L2 Instruction and L2 Data}

\begin{table}[H]
    \centering
    \begin{tabular}{llrrrrr}
        \toprule
        Trace & Policy & IPC & Speedup & L2C load MPKI & L2 instr.\ MPKI & L2 data MPKI \\
        \midrule
        \todo{} & & & & & & \\
        \bottomrule
    \end{tabular}
    \caption{Task 2B results.}
\end{table}

\todo{plots}

% ================================================================= TASK 3
\section{Task 3: Incorporating SPP with Dynamic I-RRIP}
\label{sec:task3}

\subsection{Design}
\todo{prefetch awareness in \texttt{replacement/dyn\_irrip.l2c\_repl}}

\subsection{Rationale for handling prefetch packets}
\todo{why the replacement policy treats prefetch packets this way}

\subsection{Results}
\todo{IPC, speedup, prefetch accuracy, coverage, Load MPKI of L1D, L2C, L2 Instruction and L2 Data}

\begin{table}[H]
    \centering
    \begin{tabular}{llrrrrrrrr}
        \toprule
        Trace & Config & IPC & Speedup & Acc.\ (\%) & Cov.\ (\%) & L1D MPKI & L2C MPKI & L2 instr.\ MPKI & L2 data MPKI \\
        \midrule
        \todo{} & & & & & & & & & \\
        \bottomrule
    \end{tabular}
    \caption{Task 3 results: SPP + Dynamic I-RRIP versus SPP alone and Dynamic I-RRIP alone.}
\end{table}

\todo{plots}

\subsection{Discussion}
\todo{comparison against standalone SPP and standalone Dynamic I-RRIP}

\end{document}
